\section{Self-Attention y Transformer}

\subsection{Los cuatro pasos de Self-Attention}

Self-Attention (autoatención) es una forma especial de Attention: \textbf{Query, Key y Value provienen todos de la misma secuencia de entrada}, es decir, la secuencia calcula atención sobre sus propias posiciones.

\begin{figure}[H]
\centering
\includegraphics[width=0.85\textwidth]{slides-images/slide-70.jpg}
\caption{Los cuatro pasos de Self-Attention\protect\footnotemark}
\end{figure}
\footnotetext{Diapositivas, página 70.}

El proceso completo de cálculo de Self-Attention se divide en cuatro pasos:

\textbf{Paso 1: codificar la información de posición (Positional Encoding)}

Como ya no procesamos la secuencia paso a paso, sino que introducimos todos los elementos a la vez, necesitamos codificar explícitamente la información de posición:

$$\tilde{x}_i = x_i + p_i$$

donde $p_i$ es el vector de codificación posicional, que proporciona un identificador único para cada posición.

\begin{importantbox}{La necesidad de Positional Encoding}
Sin codificación posicional, Self-Attention no podría distinguir entre ``the cat sat on the mat'' y ``the mat sat on the cat'', porque la propia operación de Attention es \textbf{invariante ante permutaciones} (Permutation Invariant) del orden de la entrada. La codificación posicional rompe esta simetría y permite que el modelo perciba las relaciones de orden dentro de la secuencia.
\end{importantbox}

\textbf{Paso 2: extraer Query, Key y Value}

A partir de la entrada ya codificada posicionalmente, se generan las matrices Q, K y V mediante tres capas de transformación lineal distintas:

$$Q = \tilde{X} \cdot W_Q, \quad K = \tilde{X} \cdot W_K, \quad V = \tilde{X} \cdot W_V$$

\begin{itemize}
    \item $W_Q, W_K, W_V$ son tres matrices de pesos que se pueden aprender
    \item la misma entrada $\tilde{X}$, al pasar por transformaciones lineales distintas, captura aspectos diferentes de la información
\end{itemize}

\textbf{Paso 3: calcular los pesos de Attention}

Se usa el \textbf{producto punto escalado} (Scaled Dot-Product) para calcular la similitud entre Query y Key:

$$\text{Attention Score} = \text{softmax}\left(\frac{Q \cdot K^T}{\sqrt{d_k}}\right)$$

\begin{itemize}
    \item $Q \cdot K^T$: el producto punto de Query y Key, que mide la similitud
    \item $\sqrt{d_k}$: el factor de escala ($d_k$ es la dimensión de Key), que evita que valores de producto punto demasiado grandes provoquen el desvanecimiento del gradiente de softmax
    \item $\text{softmax}$: normaliza la similitud a una distribución de probabilidad
\end{itemize}

\begin{figure}[H]
\centering
\includegraphics[width=0.65\textwidth]{slides-images/slide-73.jpg}
\caption{Cálculo de la similitud mediante el producto punto escalado\protect\footnotemark}
\end{figure}
\footnotetext{Diapositivas, página 73.}

\textbf{Paso 4: extraer las características de mayor atención}

Se pondera Value con los pesos de Attention y se suma para obtener la salida final:

$$\text{Output} = \text{softmax}\left(\frac{Q \cdot K^T}{\sqrt{d_k}}\right) \cdot V$$

\begin{importantbox}{La fórmula completa de Self-Attention}
$$\text{Attention}(Q, K, V) = \text{softmax}\left(\frac{Q K^T}{\sqrt{d_k}}\right) V$$

Esta fórmula se llama Scaled Dot-Product Attention y es la unidad de cálculo central de la arquitectura Transformer.
\end{importantbox}

\subsection{Interpretabilidad de los pesos de Attention}

La matriz de pesos de Attention ofrece una forma intuitiva de interpretabilidad: cada elemento $a_{ij}$ de la matriz representa el ``grado de atención'' que la posición $i$ presta a la posición $j$.

\begin{figure}[H]
\centering
\includegraphics[width=0.75\textwidth]{slides-images/slide-75.jpg}
\caption{Visualización de la matriz de pesos de Attention: cuanto más oscuro el color, mayor el grado de atención\protect\footnotemark}
\end{figure}
\footnotetext{Diapositivas, página 75.}

Por ejemplo, en la oración ``He tossed the tennis ball to serve'', la palabra ``ball'' puede tener un peso de atención alto hacia ``tennis'' y ``tossed'', porque estas palabras son las más relevantes para entender el significado de ``ball''.

\subsection{De Attention Head a Transformer}

\begin{figure}[H]
\centering
\includegraphics[width=0.65\textwidth]{slides-images/slide-78.jpg}
\caption{Self-Attention es el bloque de construcción básico de la arquitectura Transformer\protect\footnotemark}
\end{figure}
\footnotetext{Diapositivas, página 78.}

Un cálculo completo de Self-Attention (desde la codificación posicional, pasando por la extracción de Q/K/V, hasta la salida ponderada) constituye un \textbf{Attention Head}.

\subsection{Multi-Head Attention}

Un Transformer real no usa un solo Attention Head, sino \textbf{varios} Attention Head en paralelo: esto es \textbf{Multi-Head Attention}:

$$\text{MultiHead}(Q, K, V) = \text{Concat}(\text{head}_1, \ldots, \text{head}_h) W^O$$

donde cada head usa $W_Q^i, W_K^i, W_V^i$ distintos:

$$\text{head}_i = \text{Attention}(Q W_Q^i, K W_K^i, V W_V^i)$$

\begin{importantbox}{Las ventajas de Multi-Head Attention}
La atención multicabeza permite que el modelo preste atención \textbf{simultáneamente desde distintos ángulos} a diferentes aspectos de la entrada. Por ejemplo, al procesar una oración:
\begin{itemize}
    \item un head puede concentrarse en capturar relaciones de estructura sintáctica
    \item otro head puede concentrarse en la correlación semántica
    \item un tercer head puede atender a la resolución de referencias (quién se refiere a quién)
\end{itemize}
Las salidas de los distintos head se concatenan y se fusionan mediante una transformación lineal, lo que permite al modelo aprovechar de manera integral múltiples patrones de atención.
\end{importantbox}

\subsection{Resumen de la sección}

Self-Attention, mediante cuatro pasos (codificación posicional, extracción de Q/K/V, cálculo de similitud y suma ponderada), logra modelar directamente las dependencias internas de la secuencia. No necesita procesar la secuencia paso a paso, puede calcularse en paralelo y captura directamente dependencias a cualquier distancia. Multi-Head Attention refuerza aún más la capacidad del modelo de entender la entrada desde múltiples ángulos. Estos mecanismos constituyen el núcleo de la arquitectura Transformer.

\newpage
